% Answers to Chapter 8, translated from sv/back/facit/08.tex.

\facitavsnitt{c8}{Chapter 8}{Functions}
\begin{facit}

\svar{1.1} \sv{a} A function: every $x$ gives exactly one $y$.
\sv{b} Not a function: the circle gives, for example, both $y = 4$ and $y = -4$ for $x = 0$.
\sv{c} A function: all the $x$-values are different (that $y = 4$ occurs twice is allowed).
\sv{d} Not a function: $x = 1$ gives both $y = 2$ and $y = 5$.

\svar{1.2} \sv{a} $D_f = \mathbb{R} \setminus \{3\}$, $V_f = \mathbb{R} \setminus \{0\}$
\sv{b} $D_g = [-4, \infty)$, $V_g = [0, \infty)$
\sv{c} $D_h = \mathbb{R}$, $V_h = (-\infty, 2]$

\svar{1.3} \sv{a} $k = 1.5\,\Omega/{}^{\circ}\text{C}$, $m = 100\,\Omega$
\sv{b} $137.5\,\Omega$

\svar{2.1} \sv{a} $f(10) \approx 864$: after $10$ years there are about $864$ installations.
\sv{b} After about $9$ years.

\svar{2.2} \sv{a} A power function, $P(I) = 50I^2$ (a quadratic function).
\sv{b} $P(0.5) = 12.5\,\text{W}$, $P(2) = 200\,\text{W}$
\sv{c} $D_P = [0, 3]\,\text{A}$, $V_P = [0, 450]\,\text{W}$

\svar{2.3} \sv{a} The voltage falls by $3\,\text{V}$ every second; a constant change gives a
linear relationship.
\sv{b} $U(t) = -3t + 12$
\sv{c} $t = 5\,\text{s}$. No: the voltage across a discharging capacitor does not become negative
but levels out towards $0\,\text{V}$, so the model holds only within the measured interval
$0 \leq t \leq 4\,\text{s}$.

\svar{3.1} \sv{a} $7$ \sv{b} $-11$ \sv{c} $2$ \sv{d} $11$ \sv{e} $2$ \sv{f} $-3$

\svar{3.2} \sv{a} $x = 5$ \sv{b} $x = \frac{5}{3} \approx 1.67$ \sv{c} $x = 3$ or $x = -3$
\sv{d} No: $x^2 + 2 = 0$ gives $x^2 = -2$, and no real square is negative.
\sv{e} $x = 12$

\svar{3.3} \sv{a} $\mathbb{R}$ \sv{b} $\mathbb{R} \setminus \{-2\}$ \sv{c} $[3, \infty)$
\sv{d} $\mathbb{R} \setminus \{-3, 3\}$ \sv{e} $(-\infty, 9]$

\svar{3.4} \sv{a} $[1, \infty)$ \sv{b} $(-\infty, 0]$ \sv{c} $[0, \infty)$ \sv{d} $[-3, 7]$
\sv{e} $[0, \infty)$

\svar{3.5} \sv{a} $f(x) = 3x + 4$ \sv{b} $f(x) = 3x + 2$ \sv{c} $f(x) = -2x + 3$
\sv{d} $f(x) = -3$, a constant function.

\svar{3.6} \sv{a} The slope: how much $f(x)$ changes when $x$ increases by one step.
\sv{b} The $y$-intercept, $f(0) = m$.
\sv{c} Decreasing, since $k = -4 < 0$.
\sv{d} $x = 3$
\sv{e} $m = 12\,\text{V}$ is the battery's open-circuit voltage, the terminal voltage at $I = 0$.
The slope $k = -2$ is minus the internal resistance, $R_{\text{i}} = 2\,\Omega$.

\svar{3.7} \sv{a} $f(-3) = 9$, $g(-3) = -27$
\sv{b} An even exponent removes the minus sign and an odd one keeps it: $(-x)^2 = x^2$ but
$(-x)^3 = -x^3$.
\sv{c} $V_f = [0, \infty)$, $V_g = \mathbb{R}$
\sv{d} $9$ times greater.

\svar{3.8} \sv{a} $5$, $10$ and $40$ \sv{b} $100$, $50$ and $12.5$
\sv{c} $f$ is increasing and $g$ decreasing. The base decides: $a > 1$ gives an increasing
function, $0 < a < 1$ a decreasing one.
\sv{d} $C = 5$ and $C = 100$ respectively; $C$ is the initial value $f(0)$, since $a^0 = 1$.

\svar{3.9} \sv{a} $1.05$ \sv{b} $0.80$ \sv{c} $f(n) = 2 \cdot 0.7^n\,\text{V}$
\sv{d} $0.686\,\text{V}$ \sv{e} After $2$ stages, when it is $0.98\,\text{V}$.

\svar{3.10} \sv{a} $f$: the difference between adjacent values is constant, $2$.
\sv{b} $g$: the ratio between adjacent values is constant, $2$.
\sv{c} $f(x) = 2x + 3$ \sv{d} $g(x) = 3 \cdot 2^x$ \sv{e} $f(10) = 23$, $g(10) = 3072$

\svar{3.11} \sv{a} $x = 3$ \sv{b} $x = \pm 5$ \sv{c} $x_1 = 2$, $x_2 = 3$
\sv{d} No real zeros. \sv{e} $x_1 = 0$, $x_2 = 4$

\svar{3.12} \sv{a} $0\,\text{V}$, $6\,\text{V}$ and $10.8\,\text{V}$
\sv{b} $V_U = [0, 10.8]\,\text{V}$
\sv{c} $x = 2\,\text{k}\Omega$
\sv{d} No: the denominator $1 + x$ is always greater than the numerator $x$, so
$U < 12\,\text{V}$. The voltage approaches $12\,\text{V}$ but never reaches it.

\svar{3.13} \sv{a} $10\,\text{V}$ \sv{b} $u(2) \approx 3.68\,\text{V}$,
$u(4) \approx 1.35\,\text{V}$ \sv{c} Decreasing. \sv{d} $V_u = (0, 10]\,\text{V}$
\sv{e} No: $e^{-t/2} > 0$ for all $t$, so the voltage approaches $0\,\text{V}$ but never becomes
exactly zero.

\svar{3.14} \sv{a} False: a square is never negative, so $V_f = [0, \infty)$.
\sv{b} False: $x = 0$ gives division by zero, so $D_f = \mathbb{R} \setminus \{0\}$.
\sv{c} False: a zero requires $f(x) = 0$, but $f(2) = 5$.
\sv{d} False: $x = 0$ gives both $y = 5$ and $y = -5$.
\sv{e} False: the initial value is $f(0) = 3$; $2$ is the growth factor.

\svar[Test yourself]{T} \sv{1} For example $f(x) = \sqrt{x}$ with $D_f = [0, \infty)$: the square
root of a negative number is not defined. (Or $f(x) = \frac{1}{x}$, where $x = 0$ gives division by
zero.)
\sv{2} $f(2) = 9$, $f(-1) = 6$

\end{facit}
